\documentclass{oxmathproblems}

\printanswers

\course{Understanding Analysis, 2nd Edition: Stephen Abbott}
\sheettitle{Chapter 4 - Section 4.3 - Continuous Functions} 

\begin{document}
\begin{questions}

%------------------------- Problem 1 -------------------------
\miquestion Let $g(x) = \sqrt[3]{x}$.
\begin{enumerate}[label=(\alph*)]
  \item Prove that $g$ is continuous at $c = 0$.
  \item Prove that $g$ is continuous at a point $c \neq 0$. (The identity $a^3-b^3=(a-b)(a^2+ab+b^2)$ will be helpful.)
\end{enumerate}

\begin{solution}
  \begin{enumerate}[label=(\alph*)]
    \item Pick $\delta = \epsilon^3$ then $|x| < \delta \implies |\sqrt[3]{x} - 0| < \sqrt[3]{\epsilon^3} = \epsilon$.
    \item $|\sqrt[3]{x} - \sqrt[3]{c}| = \frac{|x-c|}{|\sqrt[3]{x^2} + \sqrt[3]{xc} + \sqrt[3]{c^2}|}$. Since the denom must be positive, then
    \[\frac{|x-c|}{|\sqrt[3]{x^2} + \sqrt[3]{xc} + \sqrt[3]{c^2}|} \leq |x-c|\] Pick $\delta = \epsilon$.

    ...wait no. If the denom is less than 1, then the nume will get bigger. 
    
    Instead, since $c \neq 0$, we can always choose $\delta$ small enough so that $x$ and $c$ have the same sign (say, $\delta \leq \frac{|c|}{2}$), which would make all of the terms in the denom positive. Then:
    \[\frac{|x-c|}{|\sqrt[3]{x^2} + \sqrt[3]{xc} + \sqrt[3]{c^2}|} = \frac{|x-c|}{\sqrt[3]{x^2} + \sqrt[3]{xc} + \sqrt[3]{c^2}} \leq \frac{|x-c|}{\sqrt[3]{c^2}}\]
    Pick $\delta = \min(\frac{|c|}{2}, \epsilon\sqrt[3]{c^2})$.

    Another way is to note that $\sqrt[3]{x^2} + \sqrt[3]{xc} + \sqrt[3]{c^2}$ is a quadratic that's always positive with minimum value $\frac{3}{4}\sqrt[3]{c^2}$.
    Pick $\delta=\frac{3}{4}\epsilon\sqrt[3]{c^2}$.
  \end{enumerate}
\end{solution}

%------------------------- Problem 2 -------------------------
\miquestion In all of these, let $f$ be a function defined on all of $\textbf{R}$.
\begin{enumerate}[label=(\alph*)]
  \item Let's say that $f$ is \textit{onetinuous} at $c$ if for all $\epsilon > 0$ we can choose $\delta = 1$ and it follows that $|f(x)-f(c)| < \epsilon$
  whenever $|x-c| < \delta$. Find an example of a function that is onetinuous on all of $\textbf{R}$.
  \item Let's say that $f$ is \textit{equaltinuous} at $c$ if for all $\epsilon > 0$ we can choose $\delta = \epsilon$ and it follows that $|f(x)-f(c)| < \epsilon$
  whenever $|x-c| < \delta$. Find an example of a function that is equaltinuous on \textbf{R} that is nowhere onetinuous, or explain why there is no such function.
  \item Let's say that $f$ is \textit{lesstinuous} at $c$ if for all $\epsilon > 0$ we can choose $0 < \delta < \epsilon$ and it follows that $|f(x)-f(c)| < \epsilon$
  whenever $|x-c| < \delta$. Find an example of a function that is lesstinuous on \textbf{R} that is nowhere equaltinuous, or explain why there is no such function.
  \item Is every lesstinuous function continuous? Is every continuous function lesstinuous? Explain.
\end{enumerate}

\begin{solution}
  \begin{enumerate}[label=(\alph*)]
    \item The function $f(x) = k$ for a fixed constant $k$ is onetinuous.
    \item $f(x) = x$. It's equaltinuous but not onetinuous as given $\epsilon < 1$, $\delta = 1$ doesn't work. For example, let $\epsilon=\frac{1}{2}$ and
    $x = c+\frac{1}{2}$, then $|x-c| < 1$ but $|f(x)-f(c)| = \frac{1}{2}$ is not less than $\epsilon$.
    \item Intuitively, for $\delta < \epsilon$ to work, we need a function that grows faster than $f(x)=x$. 
    
    $f(x) = 3x$ works. Solve the $\epsilon-\delta$ to get $\delta < \frac{\epsilon}{3}$. Since $\delta < \frac{\epsilon}{3} < \epsilon$, the function is lesstinuous.

    It's nowehere equaltinuous. Let $x = c+\frac{\epsilon}{2}$. If $\delta=\epsilon$ then $(|x-c| = \frac{\epsilon}{2}) < \delta$, but $|f(x)-f(c)| = \frac{3\epsilon}{2} > \epsilon$.
  
    $f(x)=x^2$ also works. It's lesstinuous everywhere (too lazy to type the $\epsilon-\delta$ derivation). But it's not nowhere equaltinuous. Again choose
    $x=c+\frac{\epsilon}{2}$. Then $|f(x)-f(c)| = \epsilon|c| + \frac{\epsilon^2}{4}$. If this were equaltinuous, then for \textit{all} $\epsilon > 0$
    we need $|f(x)-f(c)| < \epsilon$. But for any $\epsilon > 4$ this is false.

    \item Yes to both.
    
    Every lesstinuous function is continuous. For all $\epsilon > 0$ there's a $\delta$ that works, and that's enough for continuity. It doesn't matter what the bound on $\delta$
    is, whether is lesstinuous or equaltinuous or whatever.

    Every continuous function is lesstinuous. By definition of continuity, for all $\epsilon > 0$ there's a $\delta$ such that $|x-c| < \delta \implies |f(x)-f(c)| < \epsilon$.
    If $\delta \geq \epsilon$, then a $\delta_{2} < \epsilon$ must also work, as it satisfies $|x-c| < \delta_{2} < \epsilon \leq \delta$.
  \end{enumerate}
\end{solution}

%------------------------- Problem 3 -------------------------
\miquestion
\begin{enumerate}[label=(\alph*)]
  \item Supply a proof for Theorem 4.3.9 using the $\epsilon-\delta$ characterization of continuity.
  \item Give another proof using the sequential characterization of continuity (from Theorem 4.3.2 (iii)).
\end{enumerate}

\begin{solution}
  To restate the (abridged) Theorem: we're given $f$ continuous at $c$ and $g$ continuous at $f(c)$ and we need to prove $g(f(x))$ is continuous at $c$.
  \begin{enumerate}[label=(\alph*)]
    \item To prove $g(f(x))$ continuous at $c$, then we need to prove that for all $\epsilon > 0$ there exists $\delta$ such that $|x-c| < \delta \implies |g(f(x)) - g(f(c))| < \epsilon$.
    
    From continuity of $f$ at $c$, then for all $\epsilon_{f} > 0$ there exists $\delta_{f}$ such that $|x-c| < \delta_{f} \implies |f(x)-f(c)| < \epsilon_{f}$.

    From continuity of $g$ at $f(c)$, then for all $\epsilon > 0$ there exists $\delta_{g}$ such that $|y-f(c)| < \delta_{g} \implies |g(y)-g(f(c))| < \epsilon$.

    Set $\epsilon_{f} = \delta_{g}$. Plug $f(x)$ as $y$, we have $|x-c| < \delta_{f} \implies |f(x)-f(c)| < \delta_{g} \implies |g(f(x))-g(f(c))| < \epsilon$, as required.

    Just don't get confused by the interchange between $y$ and $f(x)$. $y$ is just a placeholder name that stands for any input of $g$, of which the values of $f(x)$ is a part of.
 
    \item We need to prove that $(x_{n}) \to c \implies g(f(x_{n})) \to g(f(c))$.
    
    From continuity of $f$ at $c$, we have $(x_{n}) \to c \implies f(x_{n}) \to f(c)$.

    From continuity of $g$ at $f(c)$, we have $(y_{n}) \to f(c) \implies g(y_{n}) \to g(f(c))$.

    Plug $f(x_{n})$ as $(y_{n})$, then we have $(x_{n}) \to c \implies f(x_{n}) \to f(c) \implies g(f(x_{n})) \to g(f(c))$, as required.
  \end{enumerate}
\end{solution}


\end{questions}
\end{document}
